\documentclass{article}
\usepackage[margin=1in]{geometry}
\usepackage{longtable}
\begin{document}

\textbf{Legacy note:} This split-by-person draft complements the canonical roadmap in \texttt{Documentation/SRS.tex}.\\

\section*{Phase 7 Plan - Person 1 (Person A)}

\textbf{Source of truth:} \texttt{Documentation/SRS.tex}, Section 11.5, ``Phase 7: Advanced Research and Scale-Up''.\\
\textbf{Interpretation note:} Person 1 owns the operational scale-up track after Phase 6 stability. The focus is not new alpha logic; it is making the existing system cheaper, more reliable, easier to run for long durations, and ready to support thesis-grade experiments without data loss.\\
\textbf{Implementation note:} Person 1 should avoid premature HA complexity. Scale-up work only counts if it is grounded in the single-instance architecture that already works.

\textbf{Phase duration:} after v1 stability\\
\textbf{Phase goal:} Improve scale, reliability, and cost efficiency while preserving reproducibility.

\section*{Owned Deliverables (From SRS)}
\begin{enumerate}
\item Improve system scale, reliability, and cost efficiency.
\item Add archival compaction, cold-storage policies, and long-run dashboards.
\item Add redundancy or HA only after the single-instance architecture is proven.
\end{enumerate}

\section*{Locked Implementation Decisions}
\begin{itemize}
\item Long-run observability must come before HA.
\item Archive compaction must never destroy replayability guarantees.
\item Cold-storage policy must be explicit by data tier, not ad hoc.
\item Dashboards must serve operator questions first: lag, gaps, cost, alert volume, model shadow load.
\item Redundancy is deferred until the current single-node runtime is measured and justified.
\end{itemize}

\section*{Concrete Execution Breakdown}
\begin{longtable}{p{0.28\linewidth}p{0.32\linewidth}p{0.30\linewidth}}
\textbf{Work Item} & \textbf{Artifact} & \textbf{Done Condition} \\
Scale/reliability audit & Bottleneck and failure inventory & Main storage, replay, alerting, and ML-shadow bottlenecks are measured and prioritized \\
Archive compaction & Compaction and retention workflow & Old raw/detector partitions can be compacted or tiered without breaking auditability \\
Cold storage policy & Tiering and restore rules & Clear restore path exists for historical windows needed by research or replay \\
Long-run dashboards & Operator dashboard/report set & Lag, gaps, candidate volume, alert volume, shadow-model load, and cost can be reviewed at a glance \\
Redundancy readiness & HA decision memo & The team knows whether HA is needed, and if so which minimal design to adopt later \\
\end{longtable}

\section*{Execution Sequence (No Day Timeline)}
\textbf{Block 1: Operational profiling}
\begin{itemize}
\item Measure runtime bottlenecks for ingestion, replay, evidence, alerting, and ML shadow scoring.
\item Quantify storage growth and high-cost paths.
\item Freeze one scale-up backlog ranked by operational pain and research impact.
\end{itemize}

\textbf{Block 2: Archive compaction and cold storage}
\begin{itemize}
\item Add archive compaction policy for long-retained replay data.
\item Define hot / warm / cold retention classes for raw archives, detector input, and model artifacts.
\item Prove that one archived historical window can still be restored and replayed after compaction.
\end{itemize}

\textbf{Block 3: Long-run dashboards and service hygiene}
\begin{itemize}
\item Build operator-facing dashboards or reports for gap rate, lag, storage growth, alert volume, and shadow-score load.
\item Add routine integrity checks and daily/weekly health summaries.
\item Keep the single-instance architecture healthy before considering HA.
\end{itemize}

\textbf{Block 4: Redundancy decision}
\begin{itemize}
\item Evaluate whether HA or redundancy is warranted after actual long-run measurements.
\item If yes, produce a minimal future design rather than implementing a large premature cluster.
\item Hand Person 2 the final long-run data availability guarantees for advanced research workloads.
\end{itemize}

\section*{Handoffs to Person 2}
\begin{itemize}
\item Stable retained datasets and restore guarantees for long-horizon experiments.
\item Archive tiering rules and any compaction caveats.
\item Long-run dashboard views that help research interpret coverage and bias risk.
\item Any scale constraints that change which experiments are affordable.
\end{itemize}

\section*{Gate 7 Checklist (Person 1 Ownership)}
\begin{itemize}
\item Scale and cost bottlenecks are measured, not guessed.
\item Archive compaction and cold-storage rules exist without breaking replayability.
\item Long-run dashboard/reporting exists for operator health.
\item HA/redundancy decision is documented after proving the single-instance design first.
\end{itemize}

\end{document}
